\section{Matter in the selected vacuum}
\label{sec:matter}

We return to the vacuum $\Phi=r\Phi_0$, $r\neq0$, with $v=|\Phi|/\sqrt2=\sqrt2\,r$, and ask whether the matter of the block is compatible with
it in the sense stated at the beginning of \cref{sec:kernels}.

\begin{theorem}[Compatibility with matter]
\label{thm:compat}
\begin{enumerate}[label=(\alph*),leftmargin=2.2em]
\item \emph{Unbroken algebra.} Among the $21$ Pati--Salam generators the stabiliser of $\Phi_0$ has dimension $12$ and closes under brackets; its part
orthogonal to $\su3\oplus\su2_L$ is spanned by $t_Y=t^R_3+\tfrac16H$. Among the $23$ generators it has dimension $14$; the two further directions $J+\tfrac13H$
and $J-J\Gamma$ carry the charges $4B$ and $1-\Gamma$ and are global.
\item \emph{Vector bosons.} The matrix $M_{ab}=\langle[t_a,\Phi_0],[t_b,\Phi_0]\rangle/|\Phi_0|^2$, taken relative to the module metric, has eigenvalues $0$ (twelve
times), $4$ (eight times) and $20$ (once), so the squared vector masses are $0$, $g^2v^2$ and $5g^2v^2$. The eight carry $Y=\pm\tfrac23$ (six leptoquarks
$X$) and $\pm1$ (the two $W_R^\pm$); the one at $5g^2v^2$ is neutral, a $Z'$.
\item \emph{Couplings.} In the module trace $|t^L_3|^2=\tfrac18$ and $|t_Y|^2=\tfrac5{24}$, so
\[
\sin^2\theta_W=\frac{|t^L_3|^2}{|t^L_3|^2+|t_Y|^2}=\frac38
\]
at the matching scale, the value of the Pati--Salam and $\so{10}$ embeddings \cite{GeorgiQuinnWeinberg74}.
\item \emph{Dirac operator.} The kinetic kernels are symmetric with $K^0=\id$ and are preserved by the twelve unbroken generators, and $H(p)^2=|p|^2+y^2r^2P_\nu$.
\item \emph{Chirality.} The massless fermions are $(3,2)_{1/6}\oplus(1,2)_{-1/2}$ with $\Gamma=-1$ and $(3,1)_{2/3}\oplus(3,1)_{-1/3}\oplus(1,1)_{-1}$ with $\Gamma=+1$:
the unbroken $\su2_L$ acts on one Lorentz chirality only.
\item \emph{Anomalies.} On the block the $Y^3$, mixed gravitational, $SU(2)^2Y$, $SU(3)^2Y$ and $SU(3)^3$ anomalies vanish, with and without $\nu_R$, and there are
four left-handed $SU(2)$ doublets, an even number \cite{Witten82}.
\end{enumerate}
\end{theorem}

The formula in (c) deserves a line. With one coupling for the whole Pati--Salam algebra, the effective couplings of $t^L_3$ and $t_Y$ are in the ratio of
their norms in the module, $g'^2/g^2=|t^L_3|^2/|t_Y|^2=\tfrac{3/24}{5/24}=\tfrac35$, the familiar normalisation of hypercharge in unified models, and
$\sin^2\theta_W=g'^2/(g^2+g'^2)=\tfrac38$. Every requirement of a matter-compatible vacuum is met at this stage except mass
generation for the charged fermions, which needs the electroweak stage.

\subsection{The electroweak stage and the fermion masses}

Everything in this subsection assumes, in addition, a vacuum value of the bidoublet kernels; their dynamics and their scale are not derived. In the radiative
scenario of \cref{sec:fullK} the colour-singlet doublets have squared masses $138-4x$, $186-4x$ and the eigenvalues of $M_2$, so they sit at the Pati--Salam
scale and a light doublet requires a tuning that we do not derive.

\begin{proposition}[Electroweak stage]
\label{prop:EW}
\begin{enumerate}[label=(\roman*),leftmargin=2.2em]
\item The kernels commuting with $\su3$ and with $Q=T^L_3+Y$ form a $12$-dimensional space: two each in $(\mathbf{10},\mathbf1,\mathbf3)$ and $(\mathbf{10},\mathbf3,\mathbf1)$ and four
each in $(\mathbf{15},\mathbf2,\mathbf2)$ and $(\mathbf1,\mathbf2,\mathbf2)$. In each bidoublet the up-type kernels (no $d,e$ entries) and the down-type kernels (no $u,\nu$ entries)
form two-dimensional spaces.
\item On the charge sectors $u$ ($Q=\tfrac23$), $d$ ($-\tfrac13$), $\nu$ ($0$) and $e$ ($-1$), the $(\mathbf1,\mathbf2,\mathbf2)$ kernels give $m_u=m_\nu^D$ and $m_d=m_e$; the
$(\mathbf{15},\mathbf2,\mathbf2)$ kernels give lepton masses three times the quark masses, with relative sign $-3$ once aligned with a $(\mathbf1,\mathbf2,\mathbf2)$ kernel. For
$\Phi=\kappa_1K_1+\kappa_{15}K_{15}$ of each type,
\[
m_{\rm quark}=|\kappa_1+\kappa_{15}|,\qquad m_{\rm lepton}=|\kappa_1-3\kappa_{15}|,
\]
the Georgi--Jarlskog factor \cite{GeorgiJarlskog79} used in fits of fermion masses with $\mathbf{10}\oplus\mathbf{126}$ \cite{BabuMohapatra93}.
\item With $\Phi=v_R\Phi_0+\kappa K_{1u}$ the neutral sector has singular values $\tfrac12\bigl(\sqrt{v_R^2+4\kappa^2}\mp v_R\bigr)$, each four times: the seesaw pair
\cite{Minkowski77,MohapatraSenjanovic80}, with light mass $\approx\kappa^2/v_R$.
\item For $\Phi=v_R\Phi_0+\kappa_uK_{1u}+\kappa_dK_{1d}$ the stabiliser among the $21$ generators is nine-dimensional, $\su3\oplus\uu1_Q$, and $Q$ annihilates $\Phi$.
\end{enumerate}
\end{proposition}

\begin{example}[An electroweak spectrum]
\label{ex:EWspectrum}
At $(v_R,\kappa_u,\kappa_d)=(20,1,\tfrac12)$ the squared vector masses, for $g=1$ in the units of \cref{thm:compat}(b), are $0$ (nine times), $79.360$ (twice, $W^\pm$),
$127.872$ (once, $Z$), $6400$ (six times), $6480.64$ (twice) and $32032.1$ (once). The ratio $m_W^2/m_Z^2=0.62062$ is close to $\cos^2\theta_W=\tfrac58$, and at
$v_R=200$ it differs from $\tfrac58$ by $-4.4\times10^{-5}$: $\rho=1$ at tree level up to corrections of order $\kappa^2/v_R^2$. In the same vacuum at
$(v_R,\kappa_u,\kappa_d)=(20,3,1)$ the fermion mass eigenspaces are a light Majorana neutrino of mass $(\sqrt{436}-20)/2\approx0.440$ and a heavy one of mass
$(\sqrt{436}+20)/2\approx20.440$, each of real dimension four, the $d$ quarks and the electron with $m_d=m_e=1$, and the $u$ quarks with mass $3$. The
chirality content of the neutral pair is $\langle\Gamma\rangle=\mp\cos2\vartheta$ with $\tan2\vartheta=2\kappa_u/v_R$, which is their mixing, and the Majorana-type part of the full
mass operator is exactly $v_R\Phi_0$: $\nu_L$ acquires its Majorana mass only through the mixing.
\end{example}

What is not obtained here is equally definite: the values of $\kappa_u$ and $\kappa_d$, the potential that would select them together with $v_R$, and the
hierarchy $\kappa\ll v_R$, which in Pati--Salam and left--right models is a tuning of the scalar potential \cite{Gildener76}.

\subsection{What the vacuum does not supply}

The statements of this subsection are impossibility results within explicitly stated algebras of operators. They say where an input enters, not that no
mechanism outside those algebras could replace it.

\begin{proposition}[Chirality is field content]
\label{prop:mirror}
Let $\mathcal A$ be the algebra of operators on $S_3$ generated by the $91$ internal bivectors and $J$. Then $\gamma^0$ commutes with every element of $\mathcal A$ and
with the rotations, anticommutes with $\Gamma$ and the boosts, and maps $V_*$ isometrically onto the orthogonal rank-$64$ subspace $\gamma^0V_*$, on which $\mathcal A$
acts by the same matrices and $\Gamma$ has the opposite sign. Hence every eigenspace of an element of $\mathcal A$ has a mirror image with the same eigenvalue and the
opposite chirality, and no vacuum value of a field with values in $\mathcal A$, nor any function of such values, separates $V_*$ from $\gamma^0V_*$.
\end{proposition}

The block is selected by $X=\Gamma\chi_W$, which contains the spacetime orientation through $\Gamma=J\gamma^0\gamma^1\gamma^2\gamma^3$ and is not in $\mathcal A$. The kernel field
is defined on $V_*$ and presupposes it, so it does not escape the proposition either. The Standard-Model vacuum of \cref{thm:select0} is chiral because the
matter on which it acts already is; this is the arena's form of the fact that a vector-like gauge theory does not generate chirality through its vacuum
\cite{VafaWitten84b}.

\begin{proposition}[Hermitian order parameters keep rank five]
\label{prop:rankfive}
On the spin-up part of the block, the operators commuting with $J$, $\Gamma$ and the twelve Standard-Model generators form a twelve-dimensional algebra whose
symmetric part is six-dimensional: the commutative algebra of the six multiplet projectors $(Q,L,u,d,e,\nu)$, which are pairwise inequivalent. A Hermitian vacuum
value in it commutes with $T^R_3$ and $\BmL$, so its stabiliser contains $\su3\oplus\su2_L\oplus\uu1_{T^R_3}\oplus\uu1_{\BmL}$, of rank five; the Standard-Model algebra, of
rank four, is not the unbroken algebra of such an order parameter.
\end{proposition}

This is the block version of the statement that order parameters of adjoint type do not lower the rank \cite{Li74}. The kernel field escapes it because a
Grassmann bilinear carries charge: $[t^R_3,\Phi_0]=J\Phi_0\neq0$.

\begin{proposition}[Gauge alignment does not see chirality]
\label{prop:alignment}
Consider rank-$64$ projectors $P$ commuting with Lorentz and $J$, with the alignment energy $V(P)=\tfrac1{512}\sum_T\|[T,P]\|^2$ over the $43$ compact internal
bivectors. For $P_*$, $\Pfam$, the colour-singlet lepton projector $P_{\rm lep}$ and the vector-like $Q_v=P_L+\gamma^0P_L\gamma^0$ with $P_L=P_*\tfrac12(\id-\Gamma)$:
\begin{center}
\small
\begin{tabular}{@{}lcccc@{}}
\toprule
 & $P_*$ & $\Pfam$ & $P_{\rm lep}$ & $Q_v$\\
\midrule
$V(P)$ & $5$ & $9$ & $3$ & $5$\\
stabiliser among the $43$ & $23$ & $9$ & $37$ & $25$\\
axial, $P\gamma^0P=0$ & yes & --- & --- & no\\
\bottomrule
\end{tabular}
\end{center}
Alignment prefers $P_{\rm lep}$ to $P_*$, and $V(P_*)=V(P_L)+V(P_R)$ with $V(P_L)=V(P_R)=\tfrac52$ while $V(Q_v)=2V(P_L)=V(P_*)$: the alignment energy does not see
chirality.
\end{proposition}

Vacuum alignment \cite{Peskin80,Preskill81} selects among projectors only by the amount of gauge symmetry they keep; it neither selects the Pati--Salam block nor
distinguishes a chiral block from a vector-like one of the same rank. Finally, with the single Yukawa coupling of \cref{ass:field}, matter opposes the
Standard-Model vacuum: on $(\mathbf{10},\mathbf1,\mathbf3)$ every $F_k$
of \cref{thm:oneloop1013} is positive, and on all of $\Kc$ the Yukawa part of the second variation is $-8x(\id-\Pi_\nu)\le0$. The fermion loop lowers the energy
along every direction that moves Majorana mass away from $\nu_R$ and raises it along none; on its own it would select configurations in which the mass is
spread, which break colour. The matter sector is not the agent that selects the Standard-Model orbit. At one loop the gauge fluctuations are, and they succeed
for $x<27$ locally and, numerically, for $x<x_*$.
